\documentclass[UTF8,11pt,a4paper]{ctexart}
\usepackage[a4paper,margin=2.5cm]{geometry}
\usepackage{amsmath,amssymb,amsthm,mathtools}
\usepackage{algorithm}
\usepackage{algpseudocode}
\usepackage{tikz}
\usetikzlibrary{arrows.meta,positioning,calc,fit,backgrounds}
\usepackage{booktabs}
\usepackage{enumitem}
\usepackage{hyperref}
\usepackage{xcolor}
\usepackage{caption}
\usepackage{float}
\usepackage{cleveref}
\setlength{\parindent}{2em}
\setlength{\parskip}{0.25em}
\hypersetup{colorlinks=true,linkcolor=black,citecolor=black,urlcolor=blue}
\newtheorem{lemma}{引理}
\newtheorem{theorem}{定理}
\newcommand{\Sset}{\mathcal S_t}
\newcommand{\Uset}{\mathcal U_t}
\newcommand{\Fset}{\mathcal F_t}
\newcommand{\Iset}{\mathcal I_t}
\newcommand{\PathU}{\mathcal P(\mathcal U_t)}
\title{\textbf{Dijkstra 不是贪心算法：一种从全局到局部的正确性证明}}
\author{王圣宝\\\small 杭州师范大学}
\date{}
\begin{document}
\maketitle
\vspace{-2.5em}
\begin{abstract}
Dijkstra 算法历来被归类为贪心算法。本文认为，这种分类未能准确揭示算法确定最短路径的内在机制。我们给出一种直接的、从全局到局部的正确性证明。每一轮迭代中，算法实际上从所有终点尚未解决的路径中确定一条全局最短路径。该路径自然也是通向自身终点的最短路径。更根本的是，算法利用已经确定的最短路径，在每一轮中求出一条新的最短路径。本文将这一过程理解为动态规划。我们通过三步排除证明，说明如何将完整路径空间压缩为前沿候选，如何利用松弛操作排除较长路径，以及为什么最小值提取能够确定全局最小值。因此，算法的选择来自确定性的全局比较，而非乐观的局部贪心选择。我们认为，与传统贪心分类相比，动态规划与确定性路径排除更准确地揭示了 Dijkstra 算法的本质。
\end{abstract}
\section{引言}
自 Dijkstra 于 1959 年提出这一优雅的最短路径算法以来\cite{Dijkstra1959}，它已成为计算机科学领域最著名的算法之一。该算法具有重要的实际用途，例如地图导航中的路径规划，以及通信网络中的路由计算\cite{SedgewickWayne2011}。与此同时，它长期被视为典型的贪心算法。经典教材普遍采用这一分类\cite{BrassardBratley1988,KleinbergTardos2006,CLRS2022}。Kleinberg 和 Tardos 使用 \emph{myopically}（短视地）描述逐步进行的贪心决策。Cormen 等人认为，Dijkstra 每轮选择暂定距离最小的未解决顶点，因此采用了贪心策略；随后再证明所选顶点的距离已经达到最终最优值。Sedgewick 和 Wayne 也采用选取最小距离与松弛相结合的描述，普林斯顿大学的相关课程同样将最短路径问题列入贪心算法的教学内容\cite{SedgewickWayne2011,WayneCOS423}。Roughgarden 给出了一个特别清晰的例子：先选择一个边界上的最小候选，再证明这一候选对于所选终点是最优的\cite{Roughgarden2018}。

上述证明是严谨的，而且也利用了路径首次跨越已解决集合边界的性质。不过，它们主要证明\emph{被选中终点}的最短距离已经确定。证明通常先给出选择规则，再说明该选择为何正确。这种组织方式容易使人将算法理解为：先执行局部贪心选择，再证明它是安全的。但它没有直接揭示一个更强的事实：被选中的候选实际上是\emph{所有终点尚未解决的路径}中的全局最小值。

最短路径算法的动态规划性质也受到过研究。Bellman 的最优性原理和 Lawler 的相关论述奠定了重要基础\cite{Bellman1957,Lawler1976}。Sniedovich 从动态规划角度分析 Dijkstra 算法；Misra 则提出不动点计算的解释\cite{Sniedovich2006,Misra2001}。Lew 研究贪心算法与动态规划之间的关系，但仍将 Dijkstra 归为贪心算法\cite{Lew2006}。这些研究揭示了重要的递推结构，却没有把“所有未解决终点路径上的全局最小化”作为正确性论证的中心对象。

本文从这个全局路径空间出发，证明每一轮迭代确实找到了其中的绝对最短路径。由此，该路径终点的最短路径性质自然成立。为说明这一点，我们提出三步排除证明：首先将搜索限制在前沿；其次利用最优子结构与松弛操作，为每个前沿顶点保留最佳候选；最后从剩余候选中确定全局最小值。这三种排除恰好对应 Dijkstra 算法本身的主要操作。

这种解释改变了传统论证的逻辑方向。算法并非相信某个局部上看似有利的选择，随后才获得全局结论；相反，确定性的排除过程先确立全局最小值，再将其终点标记为已解决。每一轮都利用既有最短路径推出一条新的最短路径。因此，本文主张 Dijkstra 本质上是一种动态规划算法，将其归类为贪心算法并不准确。这一判断针对的是最小值选择的数学依据，而非否定已有正确性证明的有效性。

\section{预备知识与搜索空间划分}\label{sec:prelim}
设 $G=(V,E,w)$ 是一个有限的带权有向图，源点为 $s\in V$。为使排除论证更清晰，本文统一假定各条边的权重严格为正：
\[
  w(e)>0\qquad(e\in E).
\]
路径 $P$ 的长度 $w(P)$ 为其各边权重之和。对于从 $s$ 可达的顶点 $v$，用 $\delta(s,v)$ 表示相应的最短路径距离。不可达顶点不参与下面的有限距离最小化；若不存在可达的未解决顶点，算法即终止。

\subsection{一个启发性的起步}
源点 $s$ 本身可视为第一个已解决顶点，其距离为 $0$。第二个已解决顶点则非常容易确定：它是 $s$ 的最小权重出边所指向的某个顶点。设
\[
  m=\min_{(s,v)\in E}w(s,v).
\]
这样的单边路径长度为 $m$。任何从 $s$ 出发、包含至少两条边的路径，其第一条边的权重至少为 $m$，并且还要增加至少一条正权边。因此，任意多边路径都严格长于 $m$。最小权重的单边路径，就是当前所有非空源点出发路径中的绝对最短路径。

这一事实十分简单：\emph{向一条路径增加边，必然使路径变长。} 由此产生本文的关键问题：在后续各轮中，算法是否同样从所有当前相关的路径中找出一条绝对最短路径？本文的证明给出肯定回答。

\subsection{Dijkstra 原文中的三类顶点}
Dijkstra 在其 1959 年原始论文的 Problem 2 中，已经将顶点分为三类\cite{Dijkstra1959}：集合 A 为最短路径已经确定的顶点；集合 B 为与 A 直接相连的其余顶点，下一个顶点从 B 中选出；集合 C 包含所有剩余顶点。本文沿用这一三区域结构，但采用更便于证明的记号。

在第 $t$ 轮迭代中，设 $\Sset$ 为\emph{已解决集合}，其中各顶点满足
\[
 d(x)=\delta(s,x)\qquad(x\in\Sset),
\]
并设 $\Uset=V\setminus\Sset$ 为\emph{未解决集合}。初始时，$\Sset=\{s\}$，且 $d(s)=0$。在 $\Uset$ 内进一步定义
\[
\Fset=\{f\in\Uset:\exists x\in\Sset\text{ 使得 }(x,f)\in E\}
\]
为\emph{前沿集合}（frontier）；设 $\Iset=\Uset\setminus\Fset$ 为\emph{不可见集合}（invisible set）。因而
\[
V=\Sset\;\dot\cup\;\Uset,\qquad
\Uset=\Fset\;\dot\cup\;\Iset.
\]
这里“不可见”仅表示相关顶点尚未通过一条从已解决区域向外的边直接暴露。图~\ref{fig:spatial} 用本文符号展示了 Dijkstra 原文已有的三个区域。对于有向图，前沿由\emph{从已解决集合指向外部}的边确定。

\begin{figure}[htbp]
\centering
\begin{tikzpicture}[
    >=Latex,
    vertex/.style={circle,draw,minimum size=7mm,inner sep=1pt},
    every node/.style={font=\small}
]
% Entire vertex space
\draw[thick,rounded corners] (-0.4,-2.75) rectangle (11.1,2.75);
\node[anchor=north west] at (-0.25,2.60) {$V$};

% Solved / unsolved containers
\draw[thick,rounded corners] (0.0,-2.10) rectangle (3.75,2.10);
\node at (1.88,1.78) {$\mathcal{S}_t$（已解决）};

\draw[thick,rounded corners] (4.15,-2.10) rectangle (10.70,2.10);
\node at (7.43,1.78) {$\mathcal{U}_t$（未解决）};

% Frontier and invisible partitions
\draw[dashed,rounded corners] (4.45,-1.52) rectangle (6.40,1.35);
\node at (5.43,1.07) {$\mathcal{F}_t$};

\draw[dashed,rounded corners] (6.80,-1.52) rectangle (10.30,1.35);
\node at (8.55,1.07) {$\mathcal{I}_t$};

% Solved vertices
\node[vertex] (s) at (0.60,0.05) {$s$};
\node[vertex] (a) at (1.70,0.88) {$a$};
\node[vertex] (b) at (1.75,-0.88) {$b$};
\node[vertex] (x) at (3.00,0.05) {$x$};

\draw[->,thick] (s) -- (a);
\draw[->,thick] (a) -- (x);
\draw[->,thick] (s) -- (b);
\draw[->,thick] (b) -- (x);

% Frontier vertices and boundary-crossing edges
\node[vertex] (f1) at (5.05,0.50) {$f_1$};
\node[vertex] (f2) at (5.70,-0.70) {$f_2$};
\draw[->,very thick] (x) -- (f1);
\draw[->,very thick] (b) -- (f2);

% Invisible vertices and deeper continuations
\node[vertex] (u1) at (7.70,0.50) {$u_1$};
\node[vertex] (u2) at (8.75,-0.45) {$u_2$};
\node[vertex] (u3) at (9.45,0.62) {$u_3$};
\draw[->,dashed] (f1) -- (u1);
\draw[->,dashed] (u1) -- (u3);
\draw[->,dashed] (f2) -- (u2);
\draw[->,dashed] (u2) -- (u3);

% Mark a continuation beyond the first unsolved vertex as excluded in Step 1
\node[font=\scriptsize,fill=white,inner sep=1pt] at (6.38,0.72) {$\times$ 较长路径};

% First-entry boundary
\draw[densely dotted,thick] (3.95,-2.28) -- (3.95,2.28);
\node[rotate=90,font=\scriptsize] at (3.89,0) {首次进入边界};

% Bottom labels
\end{tikzpicture}
\caption{第 $t$ 轮迭代中的三个顶点区域。任何从源点 $s$ 到未解决顶点的路径，都必须首先进入前沿集合 $\mathcal{F}_t$。由于边权严格为正，从该处继续延伸必定增加路径长度。}
\label{fig:spatial}
\end{figure}

\subsection{首次进入性质}
上述划分具有一个简单但关键的推论。

\begin{lemma}[首次进入引理]\label{lem:firstentry}
设 $P=(v_0=s,v_1,\ldots,v_k=u)$ 是终点为 $u\in\Uset$ 的任意路径，且 $v_j$ 是 $P$ 上第一个属于 $\Uset$ 的顶点。那么 $v_{j-1}\in\Sset$，并且 $v_j\in\Fset$。
\end{lemma}
\begin{proof}
根据 $j$ 的选取，$P$ 在 $v_j$ 以前的所有顶点均属于 $\Sset$，特别地，$v_{j-1}\in\Sset$。又因为 $(v_{j-1},v_j)\in E$ 且 $v_j\in\Uset$，由 $\Fset$ 的定义可知 $v_j\in\Fset$。
\end{proof}

因此，任何从源点进入未解决区域的路径，都首先到达某个前沿顶点。即使路径最终通向 $\Iset$ 中的深层顶点，也不能绕开这一边界。该性质是第~\ref{sec:proof} 节第一步排除的依据。

\section{从全局到局部的原理与路径排除}\label{sec:principle}
\subsection{全局路径空间}
对于每个 $u\in\Uset$，记 $\mathcal P(s,u)$ 为从 $s$ 到 $u$ 的路径集合。定义第 $t$ 轮的全局待解路径空间为
\[
  \PathU=\bigcup_{u\in\Uset}\mathcal P(s,u).
\]
当前轮的目标是确定
\[
  P_t^*\in\arg\min_{P\in\PathU}w(P).
\]
这里是在\emph{所有终点尚未解决的路径}中取最小值，而不是先假定各未解决顶点的最短路径已经求出，再对它们进行比较。

\subsection{从全局最小值推出局部最短路径}
设 $P_t^*$ 的终点为 $u^*$。因为
\[
  \mathcal P(s,u^*)\subseteq\PathU,
\]
所以任意从 $s$ 到 $u^*$ 的路径 $Q$ 均满足
\[
  w(P_t^*)\le w(Q).
\]
因此，一条路径若在全局待解路径空间中取得最小值，它必然也在自身终点对应的路径集合中取得最小值。\emph{每次找到一条全局最短路径，就解决了一个新的顶点。} 剩下的问题是：Dijkstra 如何在不枚举所有路径的情况下确定这个全局最小值？

\subsection{排除原则}
为了寻找最短路径，可以排除那些已经确定比其他有效候选更长的路径。具体而言，若 $P_1$ 仍然有效且 $w(P_1)<w(P_2)$，则 $P_2$ 不可能取得当前最小值。若存在等长候选，保留其中一个即可确定一条最短路径。

这里的排除只针对\emph{当前这一轮}。一条路径在进入第一个未解决顶点后继续延伸，可能与当前全局最小值无关；但当该顶点被解决后，其后续边又可能成为重要候选。下面的三步证明利用这一区别缩小当前搜索空间，并不会丢失未来可能需要的最短路径。

\section{三步排除正确性证明}\label{sec:proof}
固定第 $t$ 轮，假设存在至少一个可达的前沿顶点。对于 $x\in\Sset$，最短距离 $d(x)=\delta(s,x)$ 已经确定。下面证明下一次最小值提取能够得到 $\PathU$ 中的全局最小值。

\subsection{第一步：前沿排除}
任取
\[
  P=(s=v_0,v_1,\ldots,v_k=u)\in\PathU.
\]
设 $v_j$ 为 $P$ 上第一个未解决顶点。根据引理~\ref{lem:firstentry}，$v_j\in\Fset$。以 $v_j$ 为终点的前缀 $P^{\mathrm{first}}=(v_0,\ldots,v_j)$ 同样属于 $\PathU$。如果 $P$ 还要越过 $v_j$，它至少会增加一条正权边，因此
\[
  w(P^{\mathrm{first}})<w(P).
\]
可见，包含两个及以上未解决顶点的路径不可能是当前全局最短路径，因为它截至第一个前沿顶点的前缀已经更短。于是本轮只需考虑
\[
  s\rightsquigarrow x\to f,
  \qquad x\in\Sset,\quad f\in\Fset,
\]
且所有中间顶点属于已解决集合，唯一的未解决顶点是终点 $f$。这说明为什么当前只需考虑前沿顶点，无须进一步探索深层未解决路径。

\subsection{第二步：松弛排除}
固定 $f\in\Fset$。第一步之后保留的每条候选路径均有形式
\[
  P=Q(s,x)\circ(x,f),\qquad x\in\Sset.
\]
首先考察到 $x$ 的前缀。由于 $x$ 已经解决，任意前缀 $Q(s,x)$ 的长度不可能小于 $d(x)=\delta(s,x)$。若它更长，则用已知最短前缀替换后可以构成一条更短的到 $f$ 候选路径。所以，对每条跨界边 $(x,f)$，只需保留候选值
\[
  d(x)+w(x,f).
\]
这就是最优子结构所发挥的作用：松弛操作直接利用已经确定的最短距离，而不必枚举此前通向 $x$ 的各种路径。

同一个前沿顶点 $f$ 可能具有多个已解决前驱。它们产生的候选值中，最小者为
\[
  D_t(f)=\min_{\substack{x\in\Sset\\(x,f)\in E}}
  \{d(x)+w(x,f)\}.
\]
其他到 $f$ 的较长候选均可排除；如果出现等长，只需保留一个代表。因此，第二步同时包含“最优前缀替换”和“同终点最小化”，二者都由松弛操作体现。经过这一步，当前搜索由前沿值集合
\[
  \{D_t(f):f\in\Fset\}
\]
表示。

\subsection{第三步：全局最小值提取}
选取
\[
 f^*\in\arg\min_{f\in\Fset}D_t(f),
 \qquad D^*=D_t(f^*).
\]
下面证明，$D^*$ 不仅是前沿候选的最小值，而且是整个 $\PathU$ 的最小值。

任取 $P\in\PathU$，设 $f$ 是其第一个未解决顶点，$x$ 是 $f$ 在路径上的前驱，$L_x$ 是从 $s$ 到 $x$ 的前缀长度。由于 $x\in\Sset$，有 $d(x)\le L_x$。又因为 $D_t(f)$ 是所有从已解决区域跨界到 $f$ 的候选中的最小值，故
\[
 D_t(f)\le d(x)+w(x,f)\le L_x+w(x,f).
\]
路径 $P$ 在 $f$ 之后的剩余部分长度非负，因而 $w(P)\ge L_x+w(x,f)$。再由 $f^*$ 的选取知 $D_t(f)\ge D_t(f^*)$。综合得到
\[
 \begin{aligned}
 w(P)&\ge L_x+w(x,f)\\
     &\ge d(x)+w(x,f)\\
     &\ge D_t(f)\\
     &\ge D_t(f^*)=D^*.
 \end{aligned}
\]
确实存在一条长度为 $D^*$ 的路径，即一条已确定的最短前缀再接上相应的跨界边。因此
\[
  D^*=\min_{P\in\PathU} w(P).
\]

\begin{theorem}[单轮迭代的正确性]\label{thm:iteration}
若所有边权均为正，且至少有一个可达的未解决顶点，则任何使 $D_t(f)$ 在 $\Fset$ 上取得最小值的顶点 $f^*$ 都满足
\[
  D_t(f^*)=\delta(s,f^*).
\]
因此，其最短路径距离可以被最终确定。
\end{theorem}
\begin{proof}
由上述不等式可知，一条达到 $D_t(f^*)$ 的路径就是 $\PathU$ 中的全局最短路径。特别地，它必然也是所有以 $f^*$ 为终点的路径中的最短路径，因为后者是 $\PathU$ 的子集。
\end{proof}

随后，已解决集合扩大，未解决集合缩小：
\[
 \mathcal S_{t+1}=\Sset\cup\{f^*\},\qquad
 \mathcal U_{t+1}=\Uset\setminus\{f^*\}.
\]
下一轮可以重复同样的论证。

\subsection{与标准 Dijkstra 算法的对应}
\label{sec:correspondence}
三步排除并不是额外增加的算法操作，而是对标准算法已有操作的逻辑解释，如表~\ref{tab:correspondence} 所示。前沿限制对应第一步；松弛操作维护第二步得到的最优候选值；\textsc{Extract-Min} 完成第三步。

\begin{table}[htbp]
\centering
\caption{三步排除与 Dijkstra 标准操作的对应关系。}\label{tab:correspondence}
\small
\begin{tabular}{@{}p{0.25\linewidth}p{0.27\linewidth}p{0.40\linewidth}@{}}
\toprule
排除步骤 & 标准操作 & 数学作用 \\
\midrule
第一步：前沿排除 & 前沿限制 & 排除越过第一个未解决顶点后继续延伸的路径。\\
第二步：松弛排除 & 暂定距离的松弛更新 & 使用最优前缀，并为每个前沿顶点保留最佳候选。\\
第三步：全局最小值 & \textsc{Extract-Min} & 从全部剩余候选中选出全局最小值。\\
\bottomrule
\end{tabular}
\end{table}

为了更明确地说明三步排除证明与原算法流程的对应，算法~\ref{alg:standard} 给出一份简洁的 Dijkstra 伪代码，并在相应操作旁标注排除步骤。初始时，源点 $s$ 已加入已解决集合，算法立即对其出边进行松弛，因此第一次最小值提取使用的前沿值已经计算完成。在之后的每一轮中，提取之后进行的松弛，为下一次提取更新前沿候选值。所以，虽然循环内部先执行 \textsc{Extract-Min} 再执行松弛，但每次提取所依赖的状态，实际上已经完成了第二步。第一步则隐含在状态中：尚未由已解决集合暴露的未解决顶点，其距离标签仍为 $\infty$；只要还存在可达的前沿候选，它就不会被选中。

\begin{algorithm}[htbp]
\caption{以源点作为初始已解决顶点的 Dijkstra 算法}\label{alg:standard}
\begin{algorithmic}[1]
\State $d(s)\gets 0$；对 $v\ne s$，置 $d(v)\gets\infty$
\State $\mathcal S\gets\{s\}$；$Q\gets V\setminus\{s\}$（以 $d$ 为键的优先队列）
\For{每条 $(s,v)\in E$ 且 $v\in Q$}
  \State $d(v)\gets\min\{d(v),w(s,v)\}$；更新 $Q$ 中 $v$ 的键值 \Comment{第二步：初始松弛}
\EndFor
\While{$Q\ne\varnothing$}
  \State $u\gets\operatorname{ExtractMin}(Q)$ \Comment{第三步：前沿最小值}
  \If{$d(u)=\infty$} \State \textbf{break} \EndIf
  \State $\mathcal S\gets\mathcal S\cup\{u\}$
  \For{每条 $(u,v)\in E$ 且 $v\in Q$}
    \If{$d(u)+w(u,v)<d(v)$}
      \State $d(v)\gets d(u)+w(u,v)$；更新 $Q$ 中 $v$ 的键值 \Comment{第二步：松弛}
    \EndIf
  \EndFor
\EndWhile
\end{algorithmic}
\end{algorithm}

Dijkstra 原文已经采用三个顶点集合，并反复执行候选更新与最小值选择\cite{Dijkstra1959}。本文证明揭示了其背后的全局路径最小化依据，解释的是原有算法，并非另造一套算法。

\section{讨论：动态规划而非贪心选择}\label{sec:discussion}
\subsection{动态规划递推}
递推结构直接来自第二步。对任意前沿顶点 $f$，
\[
  D_t(f)=\min_{\substack{x\in\Sset\\(x,f)\in E}}
  \{d(x)+w(x,f)\}.
\]
右端所有 $d(x)$ 都是已经确定的最短距离。前沿全局最小值确定一条新的最短路径，它又可以用于后续迭代。这样，已解决的最短路径集合通过不断利用最优子结构逐步扩大。随着算法不断缩小未解决终点集合，先后确定的最短距离形成非递减序列。

这一递推可以与熟悉的 Fibonacci 数列计算比较：
\[
 F_0=0,\qquad F_1=1,\qquad F_n=F_{n-1}+F_{n-2}.
\]
两者都利用已经解决的子问题，构造新的准确结果。但依赖结构不同：Fibonacci 的下一个下标事先固定；Dijkstra 的下一个顶点由当前前沿及其边权决定。尽管如此，两者都不需要猜测下一个已解决子问题的数值。这种逐步递推的共同结构，支持将 Dijkstra 理解为动态规划算法。

\subsection{为什么贪心分类不正确}
标准贪心解释强调即时的最小暂定距离选择。在通常理解中，这样的选择依据当前可见的局部准则，而不会预先检查全部未来路径。Kleinberg 和 Tardos 明确使用 \emph{myopic}（短视的）描述贪心算法的逐步决策\cite{KleinbergTardos2006}。这一描述容易使人认为，算法先作出一个乐观的局部选择，再通过后续证明说明它确实安全。

教材中的正确性证明通常也采用这样的顺序：先选择最小暂定标签，再证明它对于被选顶点已是最终最短距离。这些证明在数学上是正确的，也比较了跨越已解决边界的路径。但其直接目标是特定终点的最优性，而不是首先对\emph{所有未解决终点路径的并集}求最小值。本文的全局到局部证明把后一种更强的最小化置于论证起点。

我们的三步排除在将下一顶点最终确定之前，便证明当前全局路径空间中没有任何路径短于被选边界候选。因此，这一最小值不是乐观的、短视的或盲目的局部猜测，而是完整、定向排除所有较长路径的确定性结果。此外，算法利用已经求出的最短路径不断确定新的最短路径。因此，本文否定贪心分类，并认为动态规划才是对 Dijkstra 算法更准确的解释。最小值提取只是确定性全局排除的最后表现，而不是算法正确性的根本来源。

\section{结论}
本文给出了 Dijkstra 算法的一种从全局到局部的正确性证明。三个确定性的排除步骤，依次将所有未解决终点路径缩减为前沿候选，再压缩为每个前沿顶点对应的一个最优候选，最终确定当前全局最短路径。该路径终点的最短路径最优性，则直接由全局最小性推出。

算法每轮利用此前已经确定的最短路径，求出一条新的最短路径。这是一种动态规划过程，而不是推测性的局部贪心选择。由此，Dijkstra 与 Bellman--Ford 和 Floyd--Warshall 等早已普遍被理解为动态规划的经典最短路径算法，可以置于一个更加统一的理论框架下\cite{CLRS2022}。这为主要经典最短路径算法提供了更一致的概念分类。

\section*{致谢}
作者感谢 Zimo Wang 所提供的有益讨论，这些讨论使本文的核心思想更加清晰。作者同时感谢人工智能大语言模型 Gemini 和 ChatGPT 在稿件起草、翻译和编辑润色过程中提供的辅助。

\begin{thebibliography}{99}

\bibitem{Dijkstra1959}
E.~W. Dijkstra,
``A note on two problems in connexion with graphs,''
\emph{Numerische Mathematik}, vol.~1, pp.~269--271, 1959.
doi:10.1007/BF01386390.

\bibitem{SedgewickWayne2011}
R. Sedgewick and K. Wayne,
\emph{Algorithms}, 4th ed.
Upper Saddle River, NJ: Addison-Wesley Professional, 2011.

\bibitem{BrassardBratley1988}
G. Brassard and P. Bratley,
\emph{Algorithmics: Theory and Practice}.
Englewood Cliffs, NJ: Prentice Hall, 1988.

\bibitem{KleinbergTardos2006}
J. Kleinberg and E. Tardos,
\emph{Algorithm Design}.
Boston, MA: Pearson/Addison-Wesley, 2006.

\bibitem{CLRS2022}
T.~H. Cormen, C.~E. Leiserson, R.~L. Rivest, and C. Stein,
\emph{Introduction to Algorithms}, 4th ed.
Cambridge, MA: MIT Press, 2022.

\bibitem{WayneCOS423}
K. Wayne,
``COS 423: Analysis of Algorithms---Greedy Algorithms: Shortest Path,''
Princeton University, course materials.
[Online]. Available: \url{https://www.cs.princeton.edu/~wayne/cs423/}.

\bibitem{Roughgarden2018}
T. Roughgarden,
\emph{Algorithms Illuminated, Part 2: Graph Algorithms and Data Structures}.
San Francisco, CA: Soundlikeyourself Publishing, LLC, 2018.

\bibitem{Bellman1957}
R. Bellman,
\emph{Dynamic Programming}.
Princeton, NJ: Princeton University Press, 1957.

\bibitem{Lawler1976}
E.~L. Lawler,
\emph{Combinatorial Optimization: Networks and Matroids}.
New York, NY: Holt, Rinehart and Winston, 1976.

\bibitem{Sniedovich2006}
M. Sniedovich,
``Dijkstra's algorithm revisited: The dynamic programming connexion,''
\emph{Control and Cybernetics}, vol.~35, no.~3, pp.~599--620, 2006.

\bibitem{Misra2001}
J. Misra,
``A walk over the shortest path: Dijkstra's algorithm viewed as fixed-point computation,''
\emph{Information Processing Letters}, vol.~77, nos.~2--4, pp.~197--200, 2001.
doi:10.1016/S0020-0190(00)00202-7.

\bibitem{Lew2006}
A. Lew,
``Canonical greedy algorithms and dynamic programming,''
\emph{Control and Cybernetics}, vol.~35, no.~3, pp.~621--643, 2006.

\end{thebibliography}


\end{document}
